\section{Entregables del plan de trabajo}\label{apx:entregables}
\justifying

La Tabla~\ref{tab:entregablesReales} replica el calendario de actividades comprometido en la propuesta (198~horas en total para el Trimestre 2026-Primavera) con el estado real de cada entregable al cierre del proyecto.

\begin{longtable}{p{0.04\textwidth} p{0.34\textwidth} p{0.06\textwidth} p{0.44\textwidth}}
    \caption{Entregables prometidos en la propuesta y su estado real.}
    \label{tab:entregablesReales}\\
    \toprule
    \textbf{No.} & \textbf{Actividad} & \textbf{Hrs.} & \textbf{Estado y evidencia} \\
    \midrule
    \endfirsthead

    \toprule
    \textbf{No.} & \textbf{Actividad} & \textbf{Hrs.} & \textbf{Estado y evidencia} \\
    \midrule
    \endhead

    1 & Diseño e implementación de la base de datos & 20 & \textbf{Completo}. 21 migraciones, seis apps, PostgreSQL 16. Diagrama ER en Figura~\ref{fig:er}. \\
    \midrule
    2 & Diseño de las interfaces por módulo & 10 & \textbf{Completo}. UI implementada en React 19 + Tailwind, con componentes reutilizables (\texttt{Button}, \texttt{Alert}, \texttt{Badge}, \texttt{Table}, \texttt{Modal}, \texttt{FormField}, \emph{etc}). Capturas en Sección~\ref{sec:resultados}. \\
    \midrule
    3 & Desarrollo del módulo de autenticación & 25 & \textbf{Completo}. \texttt{accounts/} + 491 líneas de tests (\texttt{tests/test\_auth.py}). Login JWT con \emph{throttling}. \\
    \midrule
    4 & Desarrollo del módulo de cartas temáticas & 25 & \textbf{Completo}. \texttt{documentos/} (parte \texttt{CartaTematica}) + tests en \texttt{tests/test\_documentos.py}. Sin exportación PDF/DOCX (véase Sección~\ref{sec:analisis}). \\
    \midrule
    5 & Desarrollo del módulo de requisitos de recuperación & 25 & \textbf{Completo}. \texttt{documentos/} (parte \texttt{RequisitoRecuperacion}). Mismo alcance/limitaciones que cartas. \\
    \midrule
    6 & Desarrollo del módulo de autoevaluación docente & 25 & \textbf{Ampliado}. \texttt{autoevaluacion/} con \emph{form-builder}, 8 tipos de pregunta, \emph{scoring} ponderado y versionado (9 incrementos PR1--PR9). 2\,044 líneas de tests. \\
    \midrule
    7 & Desarrollo del módulo de generación de reportes & 25 & \textbf{Completo}. \texttt{reportes/} + 390 líneas de tests. Export XLSX en cliente (no PDF). \\
    \midrule
    8 & Integración del sistema con la base de datos & 10 & \textbf{Completo}. ORM operativo, \emph{fixtures}, comando \texttt{cargar\_catalogos\_csv} y \emph{scripts} \texttt{seed\_users} / \texttt{seed\_demo}. \\
    \midrule
    9 & Prueba y corrección de errores & 8 & \textbf{Completo}. $\approx$4\,186 líneas de tests + \emph{smoke test} end-to-end. Historial git con \emph{fixes} rastreables (por ejemplo \emph{fix(seguridad)}, \emph{fix(autoevaluacion)}, \emph{fix(catalogos)}). \\
    \midrule
    10 & Documentación para usuarios & 10 & \textbf{Completo}. \texttt{docs/manuales/usuario.tex} y \texttt{docs/manuales/instalacion.tex}. \\
    \midrule
    11 & Redacción del reporte final del proyecto & 15 & \textbf{Completo}. Este documento (\texttt{docs/reporte/main.pdf}) con nueve secciones, dos anexos, índices y bibliografía. \\
    \midrule
    \multicolumn{4}{l}{\emph{Entregables adicionales no contemplados originalmente:}} \\
    \midrule
    E1 & Vista pública de exploración (cascada periodo→programa→UEA) & --- & \textbf{Entregado} (extra). Beneficia a estudiantes al elegir carga. \\
    \midrule
    E2 & Catálogos administrables completos + importación CSV & --- & \textbf{Entregado} (extra). Departamentos, licenciaturas, posgrados, áreas, UEA, periodos. \\
    \midrule
    E3 & Suite de pruebas de $\sim$4\,200 líneas en backend & --- & \textbf{Entregado} (extra). Ver \texttt{backend/tests/}. \\
    \midrule
    E4 & Documentación OpenAPI/Swagger auto-generada & --- & \textbf{Entregado} (extra). Vía \emph{drf-spectacular} en \texttt{/api/schema/} y \texttt{/api/docs/}. \\
    \bottomrule
\end{longtable}

Los manuales de instalación y de usuario, que originalmente se distribuían como PDFs independientes, se integran en este mismo documento como parte de este anexo: la Subsección~\ref{apx:manual-instalacion} reproduce el \emph{Manual de instalación} y la Subsección~\ref{apx:manual-usuario} el \emph{Manual de usuario}. Adicionalmente, la carpeta digital de entrega a la Coordinación de Estudios contiene los siguientes artefactos:

\begin{itemize}
    \item \texttt{reporte\_final.pdf} --- compilación íntegra de este documento (incluye los manuales embebidos como subsecciones de este anexo), sin restricciones.
    \item \texttt{codigo\_fuente.zip} --- generado con \texttt{git archive HEAD}, descomprimible sin contraseña.
    \item \texttt{manual\_instalacion.pdf} y \texttt{manual\_usuario.pdf} --- copias sueltas de los mismos manuales incluidos como subsecciones, útiles cuando se requiera consultarlos aislados del reporte.
    \item \texttt{cobertura/} --- reporte HTML de cobertura de tests generado con \texttt{pytest --cov}.
\end{itemize}

\subsection{Manual de instalación}\label{apx:manual-instalacion}
\includepdf[pages=2-, pagecommand={\thispagestyle{fancy}}]{../manuales/instalacion/main.pdf}

\subsection{Manual de usuario}\label{apx:manual-usuario}
\includepdf[pages=2-, pagecommand={\thispagestyle{fancy}}]{../manuales/usuario/main.pdf}

\newpage
